\begin{derivation}[从\eqnrefs{eq:lienard_wiechert_radiation_field}到\eqnrefs{eq:lienard_total_power_from_angular_flux}]

从远区辐射场
\[
\mathbf E_{\rm rad}
=
\frac{q_{\mathrm e}}{4\pi\varepsilon_0cR}
\frac{
\mathbf n\times[
(\mathbf n-\boldsymbol\beta_{\rm rel})
\times\dot{\boldsymbol\beta}_{\rm rel}]
}
{\mathfrak d_{\rm ret}^{\,3}},
\qquad
\mathfrak d_{\rm ret}
=
1-\mathbf n\cdot\boldsymbol\beta_{\rm rel}
\]
出发。远区场横向于传播方向，因此
\[
\mathbf n\cdot\mathbf E_{\rm rad}=0,
\qquad
\mathbf B_{\rm rad}
=
\frac1c\mathbf n\times\mathbf E_{\rm rad}.
\]
Poynting 向量的径向分量于是为
\[
\mathbf n\cdot\mathbf S_{\rm rad}
=
\frac1{\mu_0}
\mathbf n\cdot
(\mathbf E_{\rm rad}\times\mathbf B_{\rm rad})
=
\varepsilon_0c|\mathbf E_{\rm rad}|^2.
\]
乘球面元 \(R^2\dd\Omega\)，按\emph{观察时间} \(t\) 计的角能流为
\[
\frac{\dd P_{\rm obs}}{\dd\Omega}
=
\frac{q_{\mathrm e}^2}{16\pi^2\varepsilon_0c}
\frac{
\left|
\mathbf n\times[
(\mathbf n-\boldsymbol\beta_{\rm rel})
\times\dot{\boldsymbol\beta}_{\rm rel}]
\right|^2
}
{\mathfrak d_{\rm ret}^{\,6}}.
\]

这里还不能直接把它称为“源在单位时间内发射的功率”，因为不同观察方向对应的推迟时间间隔不同。由
\[
t-t_{\rm r}
=
\frac{R(t_{\rm r})}{c}
\]
求导，并用
\[
\frac{\dd R}{\dd t_{\rm r}}
=
-\mathbf n\cdot\mathbf v
\]
得到
\[
\frac{\dd t}{\dd t_{\rm r}}
=
1-\mathbf n\cdot\boldsymbol\beta_{\rm rel}
=
\mathfrak d_{\rm ret}.
\]
因此同一辐射能量满足
\[
\frac{\dd E}{\dd t_{\rm r}\dd\Omega}
=
\frac{\dd E}{\dd t\dd\Omega}
\frac{\dd t}{\dd t_{\rm r}},
\]
从而
\[
\boxed{
\frac{\dd P_{\rm emit}}{\dd\Omega}
=
\frac{q_{\mathrm e}^2}{16\pi^2\varepsilon_0c}
\frac{
\left|
\mathbf n\times[
(\mathbf n-\boldsymbol\beta_{\rm rel})
\times\dot{\boldsymbol\beta}_{\rm rel}]
\right|^2
}
{(1-\mathbf n\cdot\boldsymbol\beta_{\rm rel})^5}.
}
\]
这就是正文\eqnrefs{eq:lienard_angular_emission_power}。分母从六次幂降到五次幂，正是推迟时间 Jacobian 的结果。

总辐射功率可由这个角分布直接积分得到。由于总功率只依赖 \(\boldsymbol\beta_{\rm rel}\) 与 \(\dot{\boldsymbol\beta}_{\rm rel}\) 的相对方向，可取
\[
\boldsymbol\beta_{\rm rel}
=
\beta_{\rm rel}\mathbf e_z,
\qquad
\dot{\boldsymbol\beta}_{\rm rel}
=
\dot\beta_\perp\mathbf e_x
+
\dot\beta_\parallel\mathbf e_z,
\]
并写
\[
u=\cos\theta,
\qquad
\dd\Omega=\dd\varphi\,\dd u.
\]
对方位角先积分。直接展开双重叉积并使用
\[
\int_0^{2\pi}\cos\varphi\,\dd\varphi=0,
\qquad
\int_0^{2\pi}\cos^2\varphi\,\dd\varphi=\pi
\]
可得
\begin{align*}
&\int_0^{2\pi}\dd\varphi\,
\left|
\mathbf n\times[
(\mathbf n-\boldsymbol\beta_{\rm rel})
\times\dot{\boldsymbol\beta}_{\rm rel}]
\right|^2\\
&\qquad=
\pi\Bigl[
2\dot\beta_\parallel^2(1-u^2)
+
\dot\beta_\perp^2
\bigl((1+\beta_{\rm rel}^2)(1+u^2)-4\beta_{\rm rel}u\bigr)
\Bigr].
\end{align*}
平行、垂直两部分因此分别需要两个一维积分：
\[
I_\parallel
=
2\pi
\int_{-1}^{1}
\frac{1-u^2}{(1-\beta_{\rm rel}u)^5}\,\dd u,
\]
\[
I_\perp
=
\pi
\int_{-1}^{1}
\frac{
(1+\beta_{\rm rel}^2)(1+u^2)-4\beta_{\rm rel}u
}
{(1-\beta_{\rm rel}u)^5}\,\dd u.
\]
这些是有理函数积分。令
\[
w=1-\beta_{\rm rel}u,
\qquad
u=\frac{1-w}{\beta_{\rm rel}},
\qquad
\dd u=-\frac{\dd w}{\beta_{\rm rel}},
\]
上下限变为 \(w=1+\beta_{\rm rel}\) 与
\(w=1-\beta_{\rm rel}\)。逐项积分 \(w^{-3},w^{-4},w^{-5}\) 后得到
\[
I_\parallel
=
\frac{8\pi}{3}
\frac1{(1-\beta_{\rm rel}^2)^3}
=
\frac{8\pi}{3}\gamma_{\rm L}^6,
\]
以及
\[
I_\perp
=
\frac{8\pi}{3}
\frac1{(1-\beta_{\rm rel}^2)^2}
=
\frac{8\pi}{3}\gamma_{\rm L}^4.
\]
代回角功率，
\begin{align*}
P_{\rm rad}
&=
\frac{q_{\mathrm e}^2}{16\pi^2\varepsilon_0c}
\left(
I_\parallel\dot\beta_\parallel^2
+
I_\perp\dot\beta_\perp^2
\right)\\
&=
\frac{q_{\mathrm e}^2}{6\pi\varepsilon_0c}
\left(
\gamma_{\rm L}^6\dot\beta_\parallel^2
+
\gamma_{\rm L}^4\dot\beta_\perp^2
\right).
\end{align*}
由于
\[
\dot{\boldsymbol\beta}_{\rm rel}
=
\frac{\mathbf a}{c},
\qquad
|\boldsymbol\beta_{\rm rel}\times\mathbf a|^2
=
\beta_{\rm rel}^2a_\perp^2,
\]
并且
\[
\gamma_{\rm L}^6
\left[
a^2-
|\boldsymbol\beta_{\rm rel}\times\mathbf a|^2
\right]
=
\gamma_{\rm L}^6a_\parallel^2
+
\gamma_{\rm L}^4a_\perp^2,
\]
最终得到
\[
\boxed{
P_{\rm rad}
=
\frac{q_{\mathrm e}^2\gamma_{\rm L}^6}
{6\pi\varepsilon_0c^3}
\left[
a^2-
|\boldsymbol\beta_{\rm rel}\times\mathbf a|^2
\right],
}
\]
即正文\eqnrefs{eq:lienard_total_power_from_angular_flux}。非相对论极限
\(\beta_{\rm rel}\to0\) 自动退化为 Larmor 功率
\[
P_{\rm rad}
\longrightarrow
\frac{q_{\mathrm e}^2a^2}{6\pi\varepsilon_0c^3}.
\]
\end{derivation}

\begin{derivation}[从\eqnrefs{eq:formation_length_general_mismatch}到\eqnrefs{eq:formation_length_ultrarelativistic}]
真空远场中 $\kappa_z=(\omega/c)\cos\theta$，而 $v=\beta_{\rm rel}c$。因此正文\eqnrefs{eq:formation_length_general_mismatch} 精确给出
\begin{equation*}
L_f
=\frac{c}{\omega\,[\beta_{\rm rel}^{-1}-\cos\theta]}.
\end{equation*}
令 $x=\gamma_{\rm L}^{-2}$，则
$\beta_{\rm rel}^{-1}=(1-x)^{-1/2}$。对函数
$f(x)=(1-x)^{-1/2}$ 在 $x=0$ 处作一阶 Taylor 展开。由于
\begin{equation*}
f''(u)=\frac{3}{4}(1-u)^{-5/2},
\end{equation*}
Lagrange 余项给出某个 $\xi\in(0,x)$ 使
\begin{equation*}
r_\beta(x)
=f(x)-1-\frac{x}{2}
=\frac{f''(\xi)}{2}x^2.
\end{equation*}
因 $f''(u)$ 在 $0\le u<x<1$ 上单调递增，
\begin{equation*}
0\le r_\beta(x)
\le\frac{3x^2}{8(1-x)^{5/2}}.
\end{equation*}

角度部分使用四阶 Taylor 公式。存在介于 $0$ 与 $\theta$ 之间的中间角 $\theta_*$，使
\begin{equation*}
\cos\theta
=1-\frac{\theta^2}{2}
+\frac{\theta^4}{24}\cos\theta_*.
\end{equation*}
所以
\begin{equation*}
r_\theta(\theta)
=1-\cos\theta-\frac{\theta^2}{2}
=-\frac{\theta^4}{24}\cos\theta_*,
\qquad
|r_\theta|\le\frac{\theta^4}{24}.
\end{equation*}
由此
\begin{align*}
\beta_{\rm rel}^{-1}-\cos\theta
&=\frac{x+\theta^2}{2}+r_\beta+r_\theta\\
&=\frac{x+\theta^2}{2}(1+\delta_f),
\end{align*}
其中
\begin{equation*}
\delta_f
=\frac{2(r_\beta+r_\theta)}{x+\theta^2}.
\end{equation*}
三角不等式给出
\begin{equation*}
|\delta_f|
\le
\boxed{
\epsilon_f
\equiv
\frac{
\dfrac{3x^2}{4(1-x)^{5/2}}
+\dfrac{\theta^4}{12}
}{x+\theta^2}}.
\end{equation*}
代回形成长度即得到正文\eqnrefs{eq:formation_length_ultrarelativistic}。若 $\epsilon_f<1$，则
\begin{align*}
\left|\frac{L_f-L_f^{(0)}}{L_f^{(0)}}\right|
&=\left|\frac1{1+\delta_f}-1\right|\\
&=\frac{|\delta_f|}{|1+\delta_f|}\\
&\le\frac{\epsilon_f}{1-\epsilon_f},
\end{align*}
其中
$L_f^{(0)}=2c/[\omega(x+\theta^2)]$。这给出了常用超相对论小角形成长度的定量适用条件。
\end{derivation}

\begin{derivation}[从\eqnrefs{eq:source_total_loss_green_power}到\eqnrefs{eq:loss_absorption_radiation_decomposition}]
先证明 Green 表达式。对 $\ee^{-\ii\omega t}$ convention，频域 Maxwell 方程给
\[
\widetilde{\mathbf E}
=\ii\mu_0\omega\,\mathbf G^R\widetilde{\mathbf J}.
\]
定义复数
\[
Z_J
\equiv
\iint\dd^3r\,\dd^3r'\,
\widetilde{\mathbf J}^{*}(\mathbf r)\cdot
\mathbf G^R(\mathbf r,\mathbf r';\omega)
\cdot\widetilde{\mathbf J}(\mathbf r').
\]
则
\[
\int\dd^3r\,\widetilde{\mathbf J}^{*}\cdot\widetilde{\mathbf E}
=\ii\mu_0\omega Z_J.
\]
电子或外部源损失给场的平均功率取负的机械做功率，
\begin{align*}
P_{\rm loss}
&=-\frac12\Re(\ii\mu_0\omega Z_J)\\
&=\frac{\mu_0\omega}{2}\Im Z_J.
\end{align*}
而
\[
\Im Z_J
=\iint\dd^3r\,\dd^3r'\,
\widetilde{\mathbf J}^{*}(\mathbf r)\cdot
\operatorname{Im}\mathbf G^R(\mathbf r,\mathbf r';\omega)
\cdot\widetilde{\mathbf J}(\mathbf r'),
\]
得到正文\eqnrefs{eq:source_total_loss_green_power}。

局域能量守恒进一步把总损耗分解为材料吸收和辐射。取非磁介质
$\mathbf B=\mu_0\mathbf H$、
$\mathbf D=\varepsilon_0\boldsymbol\varepsilon_r\mathbf E$。频域 Maxwell 方程为
\[
\nabla\times\mathbf E=\ii\omega\mu_0\mathbf H,
\qquad
\nabla\times\mathbf H=\mathbf J-\ii\omega\varepsilon_0
\boldsymbol\varepsilon_r\mathbf E.
\]
第二式取复共轭：
\[
\nabla\times\mathbf H^*
=\mathbf J^*+\ii\omega\varepsilon_0
\boldsymbol\varepsilon_r^\dagger\mathbf E^*.
\]
使用向量恒等式
\[
\nabla\cdot(\mathbf E\times\mathbf H^*)
=\mathbf H^*\cdot(\nabla\times\mathbf E)
-\mathbf E\cdot(\nabla\times\mathbf H^*)
\]
得到
\begin{align*}
\nabla\cdot(\mathbf E\times\mathbf H^*)
&=\ii\omega\mu_0|\mathbf H|^2
-\mathbf E\cdot\mathbf J^*\\
&\quad
-\ii\omega\varepsilon_0
\mathbf E\cdot\boldsymbol\varepsilon_r^\dagger\mathbf E^*.
\end{align*}
对体积 $V$ 积分并取实部。第一项纯虚，不贡献平均能流；利用
\[
\Re\!\left[-\ii\,
\mathbf E\cdot\boldsymbol\varepsilon_r^\dagger\mathbf E^*\right]
=-\mathbf E^*\cdot\boldsymbol\varepsilon_I\cdot\mathbf E,
\]
以及散度定理，得到
\begin{align*}
\frac12\Re\oint_{\partial V}\dd\mathbf A\cdot
(\mathbf E\times\mathbf H^*)
&=-\frac12\Re\int_V\dd^3r\,\mathbf J^*\cdot\mathbf E\\
&\quad
-\frac{\omega\varepsilon_0}{2}
\int_V\dd^3r\,
\mathbf E^*\cdot\boldsymbol\varepsilon_I\cdot\mathbf E.
\end{align*}
识别
\[
P_{\rm rad}
=\frac12\Re\oint\dd\mathbf A\cdot(\mathbf E\times\mathbf H^*),
\qquad
P_{\rm abs}
=\frac{\omega\varepsilon_0}{2}\int\dd^3r\,
\mathbf E^*\cdot\boldsymbol\varepsilon_I\cdot\mathbf E,
\]
以及
$P_{\rm loss}=-(1/2)\Re\int\mathbf J^*\cdot\mathbf E$，就得到
\[
\boxed{P_{\rm loss}=P_{\rm abs}+P_{\rm rad}.}
\]
对被动介质 $\boldsymbol\varepsilon_I\succeq0$，故 $P_{\rm abs}\ge0$；采用出射边界条件时 $P_{\rm rad}\ge0$。
\end{derivation}

% 数学结论及完整证明已移入正文相应公式之后。


\begin{derivation}[从\eqnrefs{eq:general_2d_conductivity_kernel}到\eqnrefs{eq:2d_sheet_anisotropic_pole_determinant}]
先从面内平移不变情形开始。令
\[
\sigma_{ab}(\boldsymbol\rho,\boldsymbol\rho';\omega)
=
\sigma_{ab}(\boldsymbol\rho-\boldsymbol\rho';\omega).
\]
采用二维 Fourier 变换
\[
f(\boldsymbol\rho)
=
\int\frac{\dd^2q}{(2\pi)^2}
f(\mathbf q)\ee^{+\ii\mathbf q\cdot\boldsymbol\rho},
\]
卷积定理把正文\eqnrefs{eq:general_2d_conductivity_kernel} 化成
\[
K_a(\mathbf q,\omega)
=
\sum_b
\sigma_{ab}(\mathbf q,\omega)
E_{\parallel,b}(\mathbf q,0,\omega).
\]

对固定 $\mathbf q$，两侧体介质中的平面波都满足无源 Maxwell 方程。先看 TM 极化。取 $\mathbf q=q\hat{\mathbf x}$，切向电场为 $E_x$、磁场为 $H_y$。由
\[
\mathbf H
=
\frac1{\mu_0\omega}\mathbf k\times\mathbf E
\]
和
$\mathbf k\cdot\varepsilon_j\mathbf E=0$，对沿 $\pm z$ 传播的 TM 波得到
\[
H_y^{(\pm)}
=
\pm
\frac{\varepsilon_0\varepsilon_j\omega}{k_{zj}}
E_x^{(\pm)}
=
\pm Y_{p,j}E_x^{(\pm)}.
\]
TE 极化同理给
\[
H_x^{(\pm)}
=
\mp
\frac{k_{zj}}{\mu_0\omega}
E_y^{(\pm)}
=
\mp Y_{s,j}E_y^{(\pm)}.
\]
所以 $Y_{p,j}$、$Y_{s,j}$ 正是把切向电场映射到切向磁场的波导纳。

现在寻找无外部入射的束缚/共振表面模。取两侧都从界面向外衰减或满足出射条件的解，并令共同切向场幅为 $E_\parallel$。切向电场连续性已经满足。磁场跳跃条件
\[
\hat{\mathbf z}\times(\mathbf H_2-\mathbf H_1)
=
\mathbf K
=
\boldsymbol\sigma\mathbf E_\parallel
\]
则给出一个齐次线性方程。对各向同性二维层，TM 与 TE 不混合，因此分别得到
\[
\bigl(Y_{p,1}+Y_{p,2}+\sigma_L\bigr)E_p=0,
\]
\[
\bigl(Y_{s,1}+Y_{s,2}+\sigma_T\bigr)E_s=0.
\]
存在非零 $E_p$ 或 $E_s$ 的条件正是正文\eqnrefs{eq:2d_sheet_surface_mode_conditions}。

若
\[
\boldsymbol\sigma
=
\begin{pmatrix}
\sigma_{pp}&\sigma_{ps}\\
\sigma_{sp}&\sigma_{ss}
\end{pmatrix}
\]
在 TE/TM 基底中含非对角元，边界方程变成
\[
\left[
\begin{pmatrix}
Y_{p,1}+Y_{p,2}&0\\
0&Y_{s,1}+Y_{s,2}
\end{pmatrix}
+
\boldsymbol\sigma
\right]
\begin{pmatrix}E_p\\E_s\end{pmatrix}
=0.
\]
非零解存在当且仅当系数矩阵行列式为零，即正文\eqnrefs{eq:2d_sheet_anisotropic_pole_determinant}。

最后说明为什么二维空间色散对电子特别重要。电子在距界面有限冲击参数处的 Fourier 电流包含连续 $q$，且近场部分在
$q>\omega/c$ 时对应倏逝真空波。远场光只在接近光锥的横向通道取样，而 EELS 的 Green 核积分会覆盖这些大 $q$ 倏逝分量；它们通过上述边界矩阵直接读取 $\sigma_L(q,\omega)$、$\sigma_T(q,\omega)$ 的非局域结构。因此二维表面电导是三维双点介电核在严格零厚度极限下的自然接口，而不是人为引入的有效参数。

最后从同一边界方程求反射振幅。对 TE 波，切向电场连续给 $t_s=1+r_s$；磁场跳跃条件给
\[
Y_{s,1}(1-r_s)-Y_{s,2}t_s=\sigma_Tt_s.
\]
消去 $t_s$ 得
\[
r_s=\frac{Y_{s,1}-Y_{s,2}-\sigma_T}{Y_{s,1}+Y_{s,2}+\sigma_T}.
\]
对 TM 波，用标准 $p$ 极化单位矢量定义入射和反射幅；反射时基矢的切向分量随传播方向约定改变，边界代数给
\[
r_p=\frac{Y_{p,2}-Y_{p,1}+\sigma_L}{Y_{p,1}+Y_{p,2}+\sigma_L}.
\]
令 $\sigma_{L,T}\to0$，两式退化为普通 Fresnel 系数；令分母为零，则恢复正文表面本征模条件。

平面平移对称使散射 Green 张量对面内坐标差作二维 Fourier 对角化。源到界面再反射回观察点的法向传播相位是
\[
\ee^{\ii k_{z1}z'}\ee^{\ii k_{z1}z}=\ee^{\ii k_{z1}(z+z')}.
\]
自由一维 Helmholtz resolvent 给 $\ii/(2k_{z1})$，再按 TE/TM 完备分解便得到正文\eqnrefs{eq:2d_sheet_reflected_green_weyl}。

对匀速平行电子，
\[
\int_{-\infty}^{\infty}\dd t\,\ee^{\ii\omega t-\ii q_xvt}
=2\pi\delta(\omega-q_xv),
\]
所以 $q_x=\omega/v$。若 $k_{z1}=\ii\kappa_1$ 且 $z=z'=b$，则
\[
\left|\ee^{\ii k_{z1}(z+z')}\right|=\ee^{-2\kappa_1b},
\]
得到正文\eqnrefs{eq:2d_eels_evanscent_impact_parameter_factor}。
\end{derivation}

% 数学结论及完整证明已移入正文相应公式之后。


\begin{derivation}[从\eqnrefs{eq:polar-optical-oscillator-eom-dp76}到\eqnrefs{eq:lattice_lst_relation}]
从正文\eqnrefs{eq:polar-optical-oscillator-eom-dp76} 的受驱极性光学模式开始：
\[
M_r(\ddot u+\Gamma_{\rm ph}\dot u+\omega_{\rm TO}^2u)=Z^*eE,
\qquad
P_{\rm ion}=n_{\rm ion}Z^*eu.
\]
采用全文的 $\ee^{-\ii\omega t}$ 约定，时间导数映射为
\[
\dot u\mapsto-\ii\omega u(\omega),
\qquad
\ddot u\mapsto-\omega^2u(\omega).
\]
因此运动方程变成代数方程
\[
M_r\left(\omega_{\rm TO}^2-\omega^2-\ii\Gamma_{\rm ph}\omega\right)u(\omega)
=Z^*eE(\omega),
\]
解得
\[
u(\omega)
=\frac{Z^*e/M_r}
{\omega_{\rm TO}^2-\omega^2-\ii\Gamma_{\rm ph}\omega}E(\omega).
\]
离子极化于是为
\begin{align*}
P_{\rm ion}(\omega)
&=n_{\rm ion}Z^*eu(\omega)\\
&=\frac{n_{\rm ion}Z^{*2}e^2/M_r}
{\omega_{\rm TO}^2-\omega^2-\ii\Gamma_{\rm ph}\omega}E(\omega).
\end{align*}

高频电子极化已经吸收到背景介电常数 $\varepsilon_\infty$ 中，所以位移场为
\[
D(\omega)=\varepsilon_0\varepsilon_\infty E(\omega)+P_{\rm ion}(\omega)
=\varepsilon_0\varepsilon(\omega)E(\omega).
\]
除以 $\varepsilon_0E$：
\[
\varepsilon(\omega)
=\varepsilon_\infty
+\frac{n_{\rm ion}Z^{*2}e^2/(\varepsilon_0M_r)}
{\omega_{\rm TO}^2-\omega^2-\ii\Gamma_{\rm ph}\omega}.
\]
通分后
\begin{align*}
\varepsilon(\omega)
=\varepsilon_\infty
\frac{\omega_{\rm TO}^2-\omega^2-\ii\Gamma_{\rm ph}\omega
+n_{\rm ion}Z^{*2}e^2/(\varepsilon_0\varepsilon_\infty M_r)}
{\omega_{\rm TO}^2-\omega^2-\ii\Gamma_{\rm ph}\omega}.
\end{align*}
定义长波纵向光学频率
\[
\omega_{\rm LO}^2
\equiv\omega_{\rm TO}^2
+\frac{n_{\rm ion}Z^{*2}e^2}
{\varepsilon_0\varepsilon_\infty M_r},
\]
分子立即化为
$\omega_{\rm LO}^2-\omega^2-\ii\Gamma_{\rm ph}\omega$，所以
\[
\varepsilon(\omega)
=\varepsilon_\infty
\frac{\omega_{\rm LO}^2-\omega^2-\ii\Gamma_{\rm ph}\omega}
{\omega_{\rm TO}^2-\omega^2-\ii\Gamma_{\rm ph}\omega},
\]
即正文\eqnrefs{eq:polar_optical_dielectric_response}。

无阻尼时，$\omega=\omega_{\rm TO}$ 使分母为零，对应横向光学共振；$\omega=\omega_{\rm LO}$ 使 $\varepsilon(\omega)=0$，对应无外加自由电荷时纵向条件
$D=\varepsilon_0\varepsilon E=0$ 的非零电场解。静态极限再取 $\omega=0$、$\Gamma_{\rm ph}=0$：
\[
\varepsilon_s\equiv\varepsilon(0)
=\varepsilon_\infty
\frac{\omega_{\rm LO}^2}{\omega_{\rm TO}^2}.
\]
整理得到
\[
\frac{\varepsilon_s}{\varepsilon_\infty}
=\frac{\omega_{\rm LO}^2}{\omega_{\rm TO}^2},
\]
即正文\eqnrefs{eq:lattice_lst_relation}。
\end{derivation}

\begin{derivation}[从\eqnrefs{eq:planar_green_weyl_definition}到\eqnrefs{eq:planar_eels_mixed_green_per_length}]
先从固定平行动量 $\boldsymbol\kappa_\parallel=(\kappa_x,\kappa_y)$ 的 $z$ 向标量 Green 核开始。介质 1 中它满足
\begin{equation*}
 (\partial_z^2+\kappa_{z1}^2)\mathscr G_z(z,z')
 =-\delta(z-z').
\end{equation*}
在 $z\neq z'$ 的区域，齐次解由推迟边界条件 $\operatorname{Im}\kappa_{z1}\ge0$ 选成
\begin{equation*}
 \mathscr G_z(z,z')
 =A\,\ee^{\ii\kappa_{z1}|z-z'|}.
\end{equation*}
为了确定 $A$，在 $z'\pm\epsilon$ 间积分微分方程：
\begin{align*}
 \int_{z'-\epsilon}^{z'+\epsilon}\dd z\,
 \partial_z^2\mathscr G_z
 +\kappa_{z1}^2
 \int_{z'-\epsilon}^{z'+\epsilon}\dd z\,\mathscr G_z
 &=-1.
\end{align*}
当 $\epsilon\to0$，第二项按 $O(\epsilon)$ 消失，所以
\begin{equation*}
 \partial_z\mathscr G_z\big|_{z'+0}
 -\partial_z\mathscr G_z\big|_{z'-0}=-1.
\end{equation*}
由
\begin{align*}
 \partial_z \ee^{\ii\kappa|z-z'|}
 &=\ii\kappa\,\operatorname{sgn}(z-z')
 \ee^{\ii\kappa|z-z'|}
\end{align*}
可得跳跃
\begin{equation*}
 2\ii\kappa_{z1}A=-1,
 \qquad
 A=\frac{\ii}{2\kappa_{z1}}.
\end{equation*}
这就是 Weyl 表示中每个平行动量通道的标量传播因子。

界面反射路径从源点 $z'$ 到界面再回到观察点 $z$，因此自由空间相位 $\ee^{\ii\kappa_{z1}|z-z'|}$ 变成
\begin{equation*}
 \ee^{\ii\kappa_{z1}(z+z')}.
\end{equation*}
对 TE、TM 两个物理偏振分别乘 Fresnel 反射系数 $r_s,r_p$，并插入对应的偏振二张量，就得到正文\eqnrefs{eq:planar_scattering_green_dyadic}。电子轨迹投影随后决定同步条件与有限长度线宽。

电子轨迹为
\begin{equation*}
 \mathbf r_e(t)=(vt,0,b),
 \qquad -\frac T2\le t\le\frac T2.
\end{equation*}
Weyl 相位含
\begin{equation*}
 \ee^{\ii\boldsymbol\kappa_\parallel\cdot
 [\mathbf R_e(t)-\mathbf R_e(t')]}
 =\ee^{\ii\kappa_xv(t-t')}.
\end{equation*}
频域 Green 泛函还含 $\ee^{\ii\omega(t'-t)}$，因此双时间积分的全部纵向相位组合成
\begin{equation*}
 \ee^{-\ii\Delta(t-t')},
 \qquad
 \Delta\equiv\omega-\kappa_xv.
\end{equation*}
双积分可先因子化：
\begin{align*}
 I_T(\Delta)
 &\equiv
 \int_{-T/2}^{T/2}\dd t\,
 \ee^{-\ii\Delta t}\\
 &=\left.
 \frac{\ee^{-\ii\Delta t}}{-\ii\Delta}
 \right|_{-T/2}^{T/2}\\
 &=\frac{2\sin(\Delta T/2)}{\Delta}\\
 &=T\,\sinc\!\left(\frac{\Delta T}{2}\right).
\end{align*}
于是
\begin{align*}
 &\int_{-T/2}^{T/2}\dd t
 \int_{-T/2}^{T/2}\dd t'\,
 \ee^{-\ii\Delta(t-t')}\\
 &\qquad=I_T(\Delta)I_T^*(\Delta)\\
 &\qquad=T^2\sinc^2\!\left(\frac{\Delta T}{2}\right).
\end{align*}
因此有限轨迹的同步线形不是 $\delta$ 函数，而是正文\eqnrefs{eq:finite-trajectory-sync-linewidth-dp83} 的 $\sinc^2$ 主峰。

第一零点来自
\begin{equation*}
 \sin\frac{\Delta T}{2}=0,
 \qquad
 \frac{|\Delta|T}{2}=\pi,
\end{equation*}
所以从中心到第一零点的角频率间隔是
\begin{equation*}
 \Delta\omega_{\rm zero}=\frac{2\pi}{T}.
\end{equation*}
半高点令
\begin{equation*}
 \sinc^2x=\frac12,
 \qquad x>0.
\end{equation*}
即求方程
\begin{equation*}
 \sin x=\frac{x}{\sqrt2}.
\end{equation*}
主峰内唯一正根为
\begin{equation*}
 x_{1/2}=1.391557\ldots.
\end{equation*}
左右半高点满足
\begin{equation*}
 \Delta_\pm=\pm\frac{2x_{1/2}}{T},
\end{equation*}
所以全高半宽
\begin{align*}
 \Delta\omega_{\rm FWHM}
 &=\Delta_+-\Delta_-\\
 &=\frac{4x_{1/2}}{T}\\
 &=\frac{5.56623\ldots}{T},
\end{align*}
即正文\eqnrefs{eq:finite-trajectory-sync-linewidth-dp83}。

再严格验证长时间极限。对任意在 $\Delta=0$ 附近连续且可积的测试函数 $f(\Delta)$，考虑
\begin{align*}
 J_T
 &\equiv\int_{-\infty}^{\infty}\dd\Delta\,
 f(\Delta)
 T\sinc^2\!\left(\frac{\Delta T}{2}\right).
\end{align*}
令 $x=\Delta T/2$，则 $\dd\Delta=2\dd x/T$：
\begin{align*}
 J_T
 &=2\int_{-\infty}^{\infty}\dd x\,
 f\!\left(\frac{2x}{T}\right)
 \sinc^2x.
\end{align*}
当 $T\to\infty$，对固定 $x$ 有 $f(2x/T)\to f(0)$，且
\begin{equation*}
 \int_{-\infty}^{\infty}\dd x\,\sinc^2x=\pi.
\end{equation*}
因此
\begin{equation*}
 \lim_{T\to\infty}J_T
 =2\pi f(0),
\end{equation*}
也就是分布恒等式
\begin{equation*}
 T\sinc^2\!\left(\frac{\Delta T}{2}\right)
 \xrightarrow[T\to\infty]{\mathcal D'}
 2\pi\delta(\Delta).
\end{equation*}

现在回到 Weyl 动量积分。长轨迹极限把
\begin{equation*}
 \delta(\omega-v\kappa_x)
\end{equation*}
插入 $\kappa_x$ 积分。由 $g(\kappa_x)=\omega-v\kappa_x$、$|g'|=v$，
\begin{equation*}
 \delta(\omega-v\kappa_x)
 =\frac1v
 \delta\!\left(\kappa_x-\frac\omega v\right).
\end{equation*}
因此
\begin{align*}
 \int_{-\infty}^{\infty}\frac{\dd\kappa_x}{2\pi}
 F(\kappa_x,\kappa_y)
 2\pi\delta(\omega-v\kappa_x)
 &=\frac1v
 F\!\left(\frac\omega v,\kappa_y\right).
\end{align*}
双时间积分提供的总因子是 $T$，而概率除以轨迹长度 $L=vT$；再乘上上式 Jacobian $1/v$ 与电流中的两个速度因子 $v^2$，有
\begin{equation*}
 \frac{T}{L}\times\frac1v\times v^2
 =1.
\end{equation*}
所以单位长度谱中不再残留额外的 $T$ 或 $v$ 前因子，$\kappa_x$ 被固定在同步值 $\omega/v$，只剩 $\kappa_y$ 积分。这正得到正文\eqnrefs{eq:planar_eels_mixed_green_per_length}。
\end{derivation}

% 数学结论及完整证明已移入正文相应公式之后。


\begin{derivation}[从\eqnrefs{eq:uniform_electron_boosted_coulomb_field}到\eqnrefs{eq:feberi_bessel_coulomb_matrix_element}]
先取给定冲击参数下的直线匀速电子轨迹，并在这一计算层级忽略环境反作用对轨迹的改变。以靶中心为原点，轨迹写成 $(\mathbf b,z)$；相应实验室系 Liénard--Wiechert 电场为

\begin{align*}
\mathbf E_{e,\perp}(\mathbf b,z)
&=-\frac{e\gamma_0}{4\pi\varepsilon_0}
\frac{\mathbf b}{(b^2+\gamma_0^2z^2)^{3/2}},\\
E_{e,z}(\mathbf b,z)
&=-\frac{e\gamma_0}{4\pi\varepsilon_0}
\frac{z}{(b^2+\gamma_0^2z^2)^{3/2}}.
\end{align*}
电偶极相互作用为 $V(z)=-\mathbf d\cdot\mathbf E_e(z)$。盒平面波矩阵元含 Fourier 因子
\[
\frac1L\int_{-\infty}^{\infty}\dd z\,
\ee^{\ii\Delta\kappa z}V(z).
\]
令 $x=\gamma_0z$，则 $\dd z=\dd x/\gamma_0$，指数变成
\[
\ee^{\ii(\Delta\kappa/\gamma_0)x}.
\]
先以 $\kappa_x$ 表示一维傅里叶波数参数，单位 ${\rm m^{-1}}$。从基本积分
\[
\int_{-\infty}^{\infty}\frac{\ee^{\ii\kappa_x x}}{\sqrt{x^2+b^2}}\dd x
=2K_0(|\kappa_x|b)
\]
出发。对 $b$ 求导：
\begin{align*}
\frac{\partial}{\partial b}
\frac1{\sqrt{x^2+b^2}}
&=-\frac{b}{(x^2+b^2)^{3/2}},\\
\frac{\partial}{\partial b}2K_0(|\kappa_x|b)
&=-2|\kappa_x|K_1(|\kappa_x|b),
\end{align*}
所以
\[
\int\dd x\,
\frac{b\ee^{\ii\kappa_x x}}{(x^2+b^2)^{3/2}}
=2|\kappa_x|K_1(|\kappa_x|b).
\]
纵向积分利用
\[
\frac{x}{(x^2+b^2)^{3/2}}
=-\frac{\partial}{\partial x}
\frac1{\sqrt{x^2+b^2}}.
\]
分部积分，边界项在 $|x|\to\infty$ 消失：
\begin{align*}
\int\dd x\,
\frac{x\ee^{\ii\kappa_x x}}{(x^2+b^2)^{3/2}}
&=-\int\dd x\,\ee^{\ii\kappa_x x}
\frac{\partial}{\partial x}
\frac1{\sqrt{x^2+b^2}}\\
&=\ii\kappa_x\int\dd x\,
\frac{\ee^{\ii\kappa_x x}}{\sqrt{x^2+b^2}}\\
&=2\ii\kappa_xK_0(|\kappa_x|b).
\end{align*}
在积分恒等式中取参数 $|\Delta\kappa|/\gamma_0$，再分别乘回横向和纵向偶极分量以及盒归一化 $1/L$，即得\eqnrefs{eq:feberi_bessel_coulomb_matrix_element}。
\end{derivation}

% 数学结论及完整证明已移入正文相应公式之后。


% 数学结论及完整证明已移入正文相应公式之后。


% 数学结论及完整证明已移入正文相应公式之后。


\begin{derivation}[从\eqnrefs{eq:far_field_source_amplitude}到\eqnrefs{eq:far_field_radiated_energy_green_si}]
真空远场的 Fourier 分量满足

\begin{equation*}
 \widetilde{\mathbf H}_{\rm far}(\omega)
 =\frac1{\mu_0c}\hat{\mathbf n}\times
 \widetilde{\mathbf E}_{\rm far}(\omega),
 \qquad
 \hat{\mathbf n}\cdot\widetilde{\mathbf E}_{\rm far}=0.
\end{equation*}
时域总辐射能为
\begin{align*}
 \frac{\dd W_{\rm rad}}{\dd\Omega}
 &=r^2\int_{-\infty}^{\infty}\dd t\,
 \hat{\mathbf n}\cdot
 [\mathbf E_{\rm far}(t)\times\mathbf H_{\rm far}(t)]\\
 &=\frac{r^2}{\mu_0c}
 \int_{-\infty}^{\infty}\dd t\,|\mathbf E_{\rm far}(t)|^2.
\end{align*}
采用
\begin{equation*}
 \mathbf E(t)=\int\frac{\dd\omega}{2\pi}\,
 \widetilde{\mathbf E}(\omega)\ee^{-\ii\omega t},
\end{equation*}
Parseval 恒等式给
\begin{equation*}
 \int\dd t\,|\mathbf E(t)|^2
 =\int\frac{\dd\omega}{2\pi}|\widetilde{\mathbf E}(\omega)|^2.
\end{equation*}
实场满足 $\widetilde{\mathbf E}(-\omega)=\widetilde{\mathbf E}^{*}(\omega)$，因此正频部分为
\begin{equation*}
 \int\dd t\,|\mathbf E(t)|^2
 =\frac1\pi\int_0^\infty\dd\omega\,
 |\widetilde{\mathbf E}(\omega)|^2.
\end{equation*}
正文的远场源--场关系是
\begin{equation*}
 \widetilde{\mathbf E}_{\rm far}
 =\ii\mu_0\omega\frac{\ee^{\ii\kappa_0r}}{r}
 \mathbf A_\infty,
 \qquad \kappa_0=\omega/c.
\end{equation*}
代入 Poynting 谱后，$r^{-2}$ 与球面面积因子相消，得到
\begin{equation*}
 \frac{\dd^2 W_{\rm rad}}{\dd\omega\dd\Omega}
 =\frac{\mu_0\omega^2}{\pi c}|\mathbf A_\infty|^2,
\end{equation*}
即正文\eqnrefs{eq:far_field_radiated_energy_green_si}。正文\eqnrefs{eq:cl_photon_probability_far_green} 随后只是把这一能量谱除以单光子能量 $\hbar\omega$ 的定义转换，不需要另行引入动力学近似。
\end{derivation}

\begin{derivation}[从\eqnrefs{eq:cl_photon_probability_far_green}到\eqnrefs{eq:free_space_current_radiation_probability}]
真空标量 Helmholtz Green 函数
$g_0(\mathbf r-\mathbf r')=\ee^{\ii\kappa_0|\mathbf r-\mathbf r'|}/[4\pi|\mathbf r-\mathbf r'|]$ 的远场系数由严格极限定义：
\begin{equation*}
\lim_{r\to\infty}
4\pi r\,\ee^{-\ii\kappa_0r}
 g_0(r\hat{\mathbf n}-\mathbf r')
=\ee^{-\ii\kappa_0\hat{\mathbf n}\cdot\mathbf r'}.
\end{equation*}
电磁 Green 张量在同一极限只保留横向投影，因此
\[
\mathbf F_{\infty,0}
(\hat{\mathbf n},\mathbf r';\omega)
=\frac1{4\pi}
(\mathbf I-\hat{\mathbf n}\hat{\mathbf n})
\ee^{-\ii\kappa_0\hat{\mathbf n}\cdot\mathbf r'}.
\]
代入 $\mathbf A_\infty$ 定义：
\begin{align*}
\mathbf A_{\infty,0}
&=\int\dd^3 r'\,
\mathbf F_{\infty,0}\cdot\widetilde{\mathbf J}(\mathbf r',\omega)\\
&=\frac1{4\pi}
(\mathbf I-\hat{\mathbf n}\hat{\mathbf n})
\int\dd^3 r'\,\widetilde{\mathbf J}(\mathbf r',\omega)
\ee^{-\ii\kappa_0\hat{\mathbf n}\cdot\mathbf r'}\\
&=\frac1{4\pi}
(\mathbf I-\hat{\mathbf n}\hat{\mathbf n})
\widetilde{\mathbf J}(\kappa_0\hat{\mathbf n},\omega).
\end{align*}
因此
\[
|\mathbf A_{\infty,0}|^2
=\frac1{16\pi^2}
\left|(\mathbf I-\hat{\mathbf n}\hat{\mathbf n})
\widetilde{\mathbf J}\right|^2.
\]
代入\eqnrefs{eq:cl_photon_probability_far_green} 后，$1/\pi$ 与 $1/(16\pi^2)$ 合并成 $1/(16\pi^3)$，得到\eqnrefs{eq:free_space_current_radiation_probability}。
\end{derivation}

\begin{derivation}[从\eqnrefs{eq:eels_source_field_green_dp68}到\eqnrefs{eq:eels_radiative_absorptive_energy_balance}]

先由外部电流对诱导场所做的负功求总损失。时域表达式为
\begin{equation*}
 W_{\rm loss}
 =-\int\dd t\dd^3 r\,
 \mathbf J(\mathbf r,t)\cdot\mathbf E(\mathbf r,t).
\end{equation*}
利用实场的正负频共轭关系，正频谱写成
\begin{equation*}
 \frac{\dd W_{\rm loss}}{\dd\omega}
 =-\frac1\pi\operatorname{Re}
 \int\dd^3 r\,
 \widetilde{\mathbf J}^{*}(\mathbf r,\omega)
 \cdot\widetilde{\mathbf E}(\mathbf r,\omega).
\end{equation*}
代入正文\eqnrefs{eq:eels_source_field_green_dp68}，
\begin{align*}
 \frac{\dd W_{\rm loss}}{\dd\omega}
 &=-\frac1\pi\operatorname{Re}
 \left[\ii\mu_0\omega
 \langle\widetilde J|G^R|\widetilde J\rangle\right]\\
 &=\frac{\mu_0\omega}{\pi}
 \langle\widetilde J|\operatorname{Im}G^R|\widetilde J\rangle,
\end{align*}
即正文\eqnrefs{eq:eels_energy_loss_green_si}。

再计算局域介质中的不可逆吸收。正文\eqnrefs{eq:material_polarization_current_dp68} 给出
\begin{equation*}
 \widetilde{\mathbf J}_{\rm mat}
 =-\ii\omega\varepsilon_0
 (\boldsymbol\varepsilon_r-\mathbf I)\widetilde{\mathbf E}.
\end{equation*}
取正频周期平均功，
\begin{align*}
 \frac{\dd W_{\rm abs}}{\dd\omega}
 &=\frac1\pi\int_{V_{\rm mat}}\dd^3 r\,
 \operatorname{Re}
 \left[\widetilde{\mathbf J}_{\rm mat}\cdot
 \widetilde{\mathbf E}^{*}\right]\\
 &=\frac{\omega\varepsilon_0}{\pi}
 \int_{V_{\rm mat}}\dd^3 r\,
 \widetilde{\mathbf E}^{*}\cdot
 \operatorname{Im}\boldsymbol\varepsilon_r
 \cdot\widetilde{\mathbf E},
\end{align*}
得到正文\eqnrefs{eq:material_absorbed_energy_spectrum_si}。

最后证明总损失如何严格分成远场辐射与材料吸收。取一个包围全部耗散材料、外边界位于无损真空的控制体 $V$。按全文 $\ee^{-\ii\omega t}$ 约定，
\begin{equation*}
 \nabla\times\mathbf E=\ii\omega\mathbf B,
 \qquad
 \nabla\times\mathbf H=\mathbf J_{\rm ext}-\ii\omega\mathbf D,
\end{equation*}
所以共轭方程为
\begin{equation*}
 \nabla\times\mathbf H^*=\mathbf J_{\rm ext}^*+\ii\omega\mathbf D^*.
\end{equation*}
代入恒等式
\begin{equation*}
 \nabla\cdot(\mathbf E\times\mathbf H^*)
 =\mathbf H^*\cdot(\nabla\times\mathbf E)
 -\mathbf E\cdot(\nabla\times\mathbf H^*)
\end{equation*}
得到复 Poynting 方程
\begin{equation*}
 \nabla\cdot(\mathbf E\times\mathbf H^*)
 =\ii\omega\mathbf H^*\cdot\mathbf B
 -\mathbf E\cdot\mathbf J_{\rm ext}^*
 -\ii\omega\mathbf E\cdot\mathbf D^*.
\end{equation*}
介质只含电损耗并取 $\mathbf B=\mu_0\mathbf H$，第一项为纯虚数，对实部没有贡献。再写
\begin{equation*}
 \mathbf D=\varepsilon_0\boldsymbol\varepsilon_r\mathbf E,
 \qquad
 X\equiv\mathbf E^*\cdot\boldsymbol\varepsilon_r\mathbf E,
\end{equation*}
则 $\mathbf E\cdot\mathbf D^*=\varepsilon_0X^*$，从而
\begin{equation*}
 \operatorname{Re}[-\ii\omega\mathbf E\cdot\mathbf D^*]
 =-\omega\varepsilon_0\operatorname{Im}X
 =-\omega\varepsilon_0
 \mathbf E^*\cdot\operatorname{Im}\boldsymbol\varepsilon_r\cdot\mathbf E.
\end{equation*}
因此
\begin{align*}
 \operatorname{Re}\nabla\cdot(\mathbf E\times\mathbf H^*)
 ={}&-\operatorname{Re}(\mathbf J_{\rm ext}^*\cdot\mathbf E)\\
 &-\omega\varepsilon_0
 \mathbf E^*\cdot\operatorname{Im}\boldsymbol\varepsilon_r\cdot\mathbf E.
\end{align*}
对 $V$ 积分并用 Gauss 定理，
\begin{align*}
 \operatorname{Re}\oint_{\partial V}\dd\mathbf S\cdot
 (\mathbf E\times\mathbf H^*)
 ={}&-\operatorname{Re}\int_V\dd^3r\,
 \mathbf J_{\rm ext}^*\cdot\mathbf E\\
 &-\omega\varepsilon_0\int_{V_{\rm mat}}\dd^3r\,
 \mathbf E^*\cdot\operatorname{Im}\boldsymbol\varepsilon_r\cdot\mathbf E.
\end{align*}
全文的正频 Fourier 归一化给每一项共同的 $1/\pi$。于是右边第一项正是前面得到的 $\dd W_{\rm loss}/\dd\omega$，第二项是 $-\dd W_{\rm abs}/\dd\omega$；把 $\partial V$ 推到无损真空远场后，左边就是 $\dd W_{\rm rad}/\dd\omega$。因此逐频有
\begin{equation*}
 \frac{\dd W_{\rm loss}}{\dd\omega}
 =\frac{\dd W_{\rm rad}}{\dd\omega}
 +\frac{\dd W_{\rm abs}}{\dd\omega},
\end{equation*}
即正文\eqnrefs{eq:eels_radiative_absorptive_energy_balance}。这里的分解来自同一个 Maxwell 能量守恒恒等式，而不是把 EELS 与 CL 当成两套独立动力学。
\end{derivation}

% 数学结论及完整证明已移入正文相应公式之后。


\begin{derivation}[从\eqnrefs{eq:pinem_time_domain_work_dp68}到\eqnrefs{eq:pinem_beta_green_unified_dp47}]
采用全文的正频率复幅约定，把真实外场写成
\[
E_{{\rm ext},z}(\mathbf r,t)
=\int_0^\infty\frac{\dd\omega}{2\pi}
\left[
E_{{\rm ext},\omega,z}(\mathbf r)\ee^{-\ii\omega t}
+E_{{\rm ext},\omega,z}^*(\mathbf r)\ee^{+\ii\omega t}
\right].
\]
沿匀速轨迹 $\mathbf r_e(t)=(\mathbf R_0,v_0t)$ 代入正文\eqnrefs{eq:pinem_time_domain_work_dp68}。对固定正频率 $\omega$，定义复功幅
\begin{align*}
\mathcal W_\omega
&\equiv
q_{\mathrm e}\int_{-\infty}^{\infty}\dd t\,
 v_0 E_{{\rm ext},\omega,z}(\mathbf R_0,v_0t)
 \ee^{-\ii\omega t}\\
&=q_{\mathrm e}\int_{-\infty}^{\infty}\dd z\,
 E_{{\rm ext},\omega,z}(\mathbf R_0,z)
 \ee^{-\ii\omega z/v_0},
\end{align*}
其中第二行只用了 $z=v_0t$ 与 $\dd z=v_0\dd t$。指数因子表明电子只抽取外场纵向 Fourier 谱在
\[
k_z=\frac{\omega}{v_0}
\]
这一同步线上对应的分量。

单个边带步长的能量是 $\hbar\omega$，因此把相干功幅除以这一能量定义无量纲生成元
\[
\beta(\omega)\equiv\frac{\mathcal W_\omega}{\hbar\omega},
\]
立刻得到正文\eqnrefs{eq:pinem_beta_green_unified_dp47}。这个除法不是经验归一化：在能量格演化算符中，$\beta$ 必须是乘在无量纲平移生成元前的无量纲系数。

若外场由外源电流产生，正文采用
\[
\mathbf E_{{\rm ext},\omega}(\mathbf r)
=\ii\mu_0\omega
\int\dd^3r'\,
\mathbf G^R(\mathbf r,\mathbf r';\omega)
\cdot\widetilde{\mathbf J}_{\rm ext}(\mathbf r',\omega).
\]
把它代回 $\beta$ 并约去 $\omega$，得到
\begin{align*}
\beta(\omega)
={}&\frac{\ii\mu_0q_{\mathrm e}}{\hbar}
\int\dd z\dd^3r'\,
 \ee^{-\ii\omega z/v_0}
\hat{\mathbf z}\cdot
\mathbf G^R((\mathbf R_0,z),\mathbf r';\omega)
\cdot\widetilde{\mathbf J}_{\rm ext}(\mathbf r',\omega).
\end{align*}
因此 PINEM 耦合与 EELS/CL 使用的是同一个推迟 Green 算符；PINEM 由外源产生相干场并沿电子轨迹读取线性振幅，EELS 则对电子自身诱导场取耗散二次型。
\end{derivation}

\begin{derivation}[从\eqnrefs{eq:planar_scattering_green_dyadic}到\eqnrefs{eq:planar_aloof_gxx_explicit}]
真空侧平面波的单位 TE 偏振定义为

\[
\hat{\mathbf s}
=\frac{\hat{\mathbf z}\times\mathbf \kappa_\parallel}{\kappa_\parallel}.
\]
若
\[
\mathbf \kappa_\parallel=\kappa_x\hat{\mathbf x}+\kappa_y\hat{\mathbf y},
\]
则
\[
\hat{\mathbf s}
=\frac{-\kappa_y\hat{\mathbf x}+\kappa_x\hat{\mathbf y}}{\kappa_\parallel},
\]
因此
\[
|\hat{\mathbf x}\cdot\hat{\mathbf s}|^2
=\frac{\kappa_y^2}{\kappa_\parallel^2}.
\]

真空 TM 偏振取
\[
\hat{\mathbf p}_1^\pm
=\frac{\kappa_\parallel\hat{\mathbf z}
\mp \kappa_{z1}\hat{\boldsymbol\kappa}_\parallel}{\kappa_0},
\qquad
\hat{\boldsymbol\kappa}_\parallel
=\frac{\mathbf \kappa_\parallel}{\kappa_\parallel}.
\]
所以
\begin{align*}
\hat{\mathbf x}\cdot\hat{\mathbf p}_1^+
&=-\frac{\kappa_{z1}}{\kappa_0}\frac{\kappa_x}{\kappa_\parallel},\\
\hat{\mathbf p}_1^-\cdot\hat{\mathbf x}
&=+\frac{\kappa_{z1}}{\kappa_0}\frac{\kappa_x}{\kappa_\parallel},
\end{align*}
从而
\[
(\hat{\mathbf x}\cdot\hat{\mathbf p}_1^+)
(\hat{\mathbf p}_1^-\cdot\hat{\mathbf x})
=-\frac{\kappa_{z1}^2}{\kappa_0^2}
\frac{\kappa_x^2}{\kappa_\parallel^2}.
\]
亚光速同步条件给 $\kappa_\parallel>\kappa_0$，故
\[
\kappa_{z1}=\ii\alpha_1,
\qquad
-\kappa_{z1}^2=\alpha_1^2.
\]
于是 TM 投影化为
\[
\frac{\alpha_1^2}{\kappa_0^2}
\frac{\kappa_x^2}{\kappa_\parallel^2}.
\]
Weyl 散射 Green 张量的公共传播因子为
\[
\frac{\ii}{2\kappa_{z1}}
\ee^{\ii \kappa_{z1}(z+z')}.
\]
取源点和观察点都在电子轨迹 $z=z'=b$：
\begin{align*}
\frac{\ii}{2\kappa_{z1}}\ee^{2\ii \kappa_{z1}b}
&=\frac{\ii}{2\ii\alpha_1}\ee^{-2\alpha_1 b}\\
&=\frac{\ee^{-2\alpha_1 b}}{2\alpha_1}.
\end{align*}
TE 投影乘该因子给
\[
\ee^{-2\alpha_1 b}\frac{\kappa_y^2}{2\alpha_1 \kappa_\parallel^2}r_s,
\]
TM 投影乘该因子给
\[
\ee^{-2\alpha_1 b}
\frac{\alpha_1 \kappa_x^2}{2\kappa_0^2\kappa_\parallel^2}r_p.
\]
两项相加得到\eqnrefs{eq:planar_aloof_gxx_explicit}。
\end{derivation}

\begin{derivation}[从\eqnrefs{eq:planar_eels_mixed_green_per_length}到\eqnrefs{eq:planar_aloof_eels_explicit_fresnel_integral}]
三维 Green 张量满足

\[
[\nabla\times\nabla\times-\kappa_0^2\varepsilon]\mathbf G^R
=\mathbf I\delta^{(3)}(\mathbf r-\mathbf r'),
\]
因此 $[\mathbf G^R]={\rm m^{-1}}$。Weyl 表示定义为
\[
\mathbf G^R
=\int\frac{\dd^2 \kappa_\parallel}{(2\pi)^2}
\ee^{\ii\mathbf \kappa_\parallel\cdot(\mathbf R-\mathbf R')}
\mathbf G_{{\rm sc},W}^R.
\]
因为 $[\dd^2 \kappa_\parallel]={\rm m^{-2}}$，必须有
\[
[\mathbf G_{{\rm sc},W}^R]={\rm m}.
\]
因此 $\dd \kappa_y\,G_{{\rm sc},W;xx}^R$ 无量纲。

前因子单位为
\begin{align*}
\left[\frac{\mu_0e^2}{\hbar}\right]
&=\frac{({\rm kg\,m\,C^{-2}})({\rm C^2})}
{{\rm J\,s}}\\
&=\frac{{\rm kg\,m}}{{\rm kg\,m^2\,s^{-1}}}\\
&={\rm s\,m^{-1}}.
\end{align*}
所以
\[
\left[\frac1L\frac{\dd P}{\dd\omega}\right]
={\rm s\,m^{-1}}.
\]
若换变量 $E_{\rm loss}=\hbar\omega$，
\[
\frac{\dd\omega}{\dd E_{\rm loss}}=\frac1\hbar,
\]
因此
\[
\frac1L\frac{\dd P}{\dd E_{\rm loss}}
=\frac1\hbar\frac1L\frac{\dd P}{\dd\omega}.
\]
\end{derivation}

\begin{derivation}[从\eqnrefs{eq:tm_surface_mode_condition}到\eqnrefs{eq:drude_surface_pole_electron_sync_frequency}]
无损 Drude 介电函数为

\[
\varepsilon_2(\omega)=1-\frac{\omega_p^2}{\omega^2}.
\]
真空--介质界面 TM 表面模色散满足
\[
\frac{\kappa_{\rm sp}}{\kappa_0}
=\sqrt{\frac{\varepsilon_2}{1+\varepsilon_2}}.
\]
在 $\kappa_y=0$ 时，电子源的同步条件是
\[
\kappa_{\rm sp}=\frac{\omega}{v}
=\frac{\kappa_0}{\beta_{\rm rel}}.
\]
平方并消去 $\kappa_0$：
\[
\frac1{\beta_{\rm rel}^2}
=\frac{\varepsilon_2}{1+\varepsilon_2}.
\]
交叉相乘：
\begin{align*}
1+\varepsilon_2
&=\beta_{\rm rel}^2\varepsilon_2,\\
1
&=(\beta_{\rm rel}^2-1)\varepsilon_2,\\
\varepsilon_2
&=-\frac1{1-\beta_{\rm rel}^2}.
\end{align*}
再代入 Drude 形式：
\begin{align*}
1-\frac{\omega_p^2}{\omega^2}
&=-\frac1{1-\beta_{\rm rel}^2},\\
\frac{\omega_p^2}{\omega^2}
&=1+\frac1{1-\beta_{\rm rel}^2}\\
&=\frac{2-\beta_{\rm rel}^2}{1-\beta_{\rm rel}^2}.
\end{align*}
所以
\[
\omega^2
=\omega_p^2\frac{1-\beta_{\rm rel}^2}{2-\beta_{\rm rel}^2}.
\]
取正频率根得到\eqnrefs{eq:drude_surface_pole_electron_sync_frequency}。
\end{derivation}
